\subsection{迭代及最终表的描述}\label{iter:subsection}

固定刚构造的一轮，包括 $D$、$G$、$q$ 和该轮的 $s$ 个分组步骤。为使总计费量节省一个 $h^a$ 的因子，取 \[
 j=\left\lceil\frac{a\log h}{\log G}\right\rceil,\qquad
 L_i=L_0G^{-i},\qquad T_i=T_0q^i \qquad(0\le i\le j).
\] 这里 $L_i,T_i$ 是完成 $i$ 轮之后的预算。还须选择 $r$ 和初始块数 $S$。较大的 $r$ 可降低只装了一部分的包的相对开销。较大的 $S$ 保证全部 $sj$ 步都有充足的组，还留下足够多的块，可在第~\ref{sec:assembly} 节中分摊公共宿主图的成本。

\paragraph{装包余量与外部宽度。} 经过一次分组，$N_b$ 除以 $r$，而 $H$ 至多增至原来的 $cr$ 倍。两者的乘积给出分配的位置实例数，其增长倍数至多为 $c$。宿主图的固定扩张可计入 $c$。因此在构造的每个阶段，都有 \begin{equation}
 H N_b\le C^{j+1}mS\log h , \label{iter:occurrences}
\end{equation} 其中 $C$ 仅依赖于单轮方案。因为 $L_0=C_0 mT_0 S\log(2+h)$，所以在以 $L_i,T_i$ 开始的轮次中，控制余量的比率满足 \[
   \frac{H N_b T_i}{rL_i}\le \frac{C_1C^{j+1}(qG)^i}{r}.
\] 取 $r$ 为不小于 $C_2(CqG)^{j+1}$ 的最小二的幂。由单轮方案固定 $C_2$，使得上述比率不超过 \eqref{eq:round-slack} 中的 $c_D$，并且 $r$ 满足每一步的下界。特别地，$r=h^{O_D(a)+o(1)}$，且与 $S$ 无关。外部宽度另有一个上界。每一轮都包含那 $ed$ 个外露比特，并且 $G^j<Gh^a$。因此对于 $i\le j$，有 \[
       \frac{ed}{L_i}=\frac{eG^i}{C_0m\log(2+h)}
           =O_D\!\left(\frac{h^{-a}}{\log h}\right).
\] 当 $h$ 足够大时，得到 $ed\le L_i/D$，恰为不可移动外部恒等站点预留的宽度。

\paragraph{副本与终止时的预算。} 任意固定 $M>0$，设 $R$ 是不小于 $h^M$ 的最小二的幂。取 $S=Rr^{sj}$，经过 $j$ 轮后恰剩 $R$ 个块。每个初始块含 $T_0$ 个副本，故 $d=T_0S$ 且 $\log d=O((\log h)^2)$。特别地，$d=n^{o(1)}$。现在 $h^a\le G^j<Gh^a$，$q\le G^{1+\delta/3}$，且 $T_0=\Theta(h^b)$。由 \eqref{eq:round-start} 和 \eqref{iter:occurrences} 得 \begin{equation}
 \begin{aligned}
 L_j&=O(md\log h\,h^{-a}),&
 T_j&=O\bigl(h^{b+(1+\delta/3)a}\bigr),\\
 HN_b&\le mS h^{O_D(a)+o(1)} .
 \end{aligned}
 \label{eq:round-terminal}
\end{equation} 位置实例数的界自始至终成立，且 $O_D(a)$ 中的常数与 $M,a,b,\beta$ 无关。由这些尺度也可知，整个过程中可以一直用同一个域。使用的尺度从 $T_0\ge4h^\beta$ 起，直至 $h$ 的一个固定次幂。选取扩张次数，使其为引理~\ref{lem:local} 中固定扩域次数的倍数，并且足够大，使域的基数在尺度 $T_j$ 上满足要求。其次数为 $O(\log h)$，故对每个 $i$ 也为 $O(\log T_i)$。

\paragraph{表的描述。} 每个位置实例携带其表的标签、各线束的宽度和这些线束在端口排序中的位置。每条关联宿主边可以用它在局部邻点列表中的位置命名，因此命名的成本可由 $D$ 控制。以 $D$ 为下标的常数可以依赖于单轮方案。每个位置实例的这些局部数据的总长度至多为 \begin{equation}
 C_3^{j+1}(\log h)^{O(1)}
   =h^{O_D(a)+o(1)} .
 \label{iter:labels}
\end{equation}

初始时，一个位置上的叠层有固定表标签，由其原表、叠层长度和各有序子列表的位置指定。这一描述的长度为 $(\log h)^{O(1)}$。在一轮之内，能够进入任何一个新标签的旧标签数，可由下述事实控制。\begin{itemize} \item 一个容量为 $DT_i$ 的聚落至多容纳 $D/\theta$ 个进入的活动站点，因为各站点的计费量至少为 $\theta T_i$。在引理~\ref{lem:local} 的局部展开中，每份输出描述使用有界个旧标签、离散选择及域元素。额外的描述长度为 $O_D(\log h)$。\item 在子轮或重置中，每个叶的负载上界，以及对同一新包有贡献的叶数上界，分别是 $T_i$ 和 $1$ 的固定倍数，倍数依赖于 $D$。每个活动核心至少贡献一个宽度为 $u\ge T_i/\operatorname{poly}(D)$ 的内部单位。单个冻结站点的计费负载即使经过固定缩放，仍有同类下界。标量站点一旦形成，随即取出。因此，每个有贡献的包都含有有界个旧标签和有界个单位。于是，单位配对列表、转子层列表以及一个端口处的有序子列表，其条目数都有界。每项至多需要 $O_D(\log h+\log r)$ 比特。这也描述了冻结站点处最终各棵树的布线。\end{itemize} 因此，若 $A$ 为当前标签长度的上界，则一步之后的上界为 $C_4A+C_4(\log h)^{C_5}$，其中包括在局部取出的标量因子的描述符。乘法常数计数的是旧标签，与扩域次数和输入字典标签的长度无关。迭代全部 $sj$ 步即证式~\eqref{iter:labels}。宿主图的全局地址只在控制整个位置时用到，无须附加在单个表的描述上。

\paragraph{总和所表示的张量。} 此时每一项由 $R$ 个保留下来的块以及在其演变过程中取出的各标量因子组成。重新分组时，外部列表与各自的站点一同置换，各乘积的每个表项都保持不变。局部张量恒等式把一个乘积换成一个和，而那 $ed$ 个已命名端口上的各个表项保持相同。着色、装包和布线在遇到平局时都按固定规则选择。唯有局部张量恒等式中的选择用来标记展开项。收缩最终的 $R$ 个块时，对最后的内部连线赋值求和。

一次这样的求值只需一部字典。一个键记录表的类型、阶段、描述和元数，并附上与表关联的比特。通过记录长度的字段或固定填充，得到唯一编码。输入字典和构造规则给每个有效键赋予域中的一个值。未用的键可任意赋值。特别地，已取出的标量的键没有关联比特。由最大元数不变量及式~\eqref{iter:labels}，最长键的长度为 \[
 O_D(T_j)+h^{O_D(a)+o(1)}
\] 推导出的表描述可能超过 $(\log h)^{O(1)}$ 比特。第~\ref{sec:assembly} 节把这段描述放进字典调用的\emph{实参}。安装的表的标签只需指定键的格式和长度。这样就分别计算了查表操作的元数，以及下一次宏操作要求输入的短标签。
